\chapter{快速起步}
\label{ch:quick-start}

第一段脚本先回答三个问题：从哪里启动，接收什么输入，产生什么结果。确认这些运行条件以后，才能把语言问题与环境问题分开排查。


\begin{figure}[htbp]
  \centering
  \begin{tikzpicture}[book mindmap,level 1 concept/.append style={minimum size=19mm,text width=15mm,font=\sffamily\bfseries\footnotesize}]
    \node[text=white,minimum size=23mm,text width=17mm,font=\sffamily\bfseries\small] {第一章\\快速起步}
      child[grow=180,level distance=33mm,concept color=DeepBlue] {
        node[text=white] {01\\输出函数}
      }
      child[grow=0,level distance=33mm,concept color=PythonYellow] {
        node {02\\项目环境}
      }
      child[grow=-90,level distance=23mm,concept color=DeepBlue!20!white] {
        node {03\\读懂报错}
      };
  \end{tikzpicture}
  \caption*{起步路线：输出、运行环境与错误定位。}
\end{figure}

\section{Python 编程第一步}

\begin{myunderstanding}
将 \code{print()} 当成小票机。要核对结果，可用它把手里的东西显示出来。
拿到 \code{return} 交回的值，还能继续用它做计算。
\end{myunderstanding}

\code{print()} 是 Python 的内置函数。默认情况下，它把传入的对象转换成文本，
用空格分隔多个对象，并在末尾换行。将下面这段保存为 \code{main.py}，
运行后先看到 \code{Hello, Python!}，下一行是 \code{5}。
签名里的 \code{a: int} 和 \code{b: int} 标注参数的预期类型，
\code{-> int} 标注返回值的预期类型。

\begin{lstlisting}[title={\ttfamily main.py}]
def add(a: int, b: int) -> int:
    return a + b  # 把结果交给调用方。


print("Hello, Python!")
print(add(2, 3))
\end{lstlisting}

\begin{pythonictip}
让 \code{add()} 只管算，把显示交给调用方。这样写测试时，可以直接比较
\code{add(2, 3)} 的结果和 \code{5}。类型注解则留给读代码的人、编辑器和静态检查工具。
\end{pythonictip}


\section{在 PyCharm 里怎么跑脚本}

将项目和解释器想成工具箱和工位。脚本都在工具箱里，选中的工位才决定用哪套工具。
新项目通常建一个 \code{.venv}，再把里面的 Python 选成项目解释器。
首次运行时，右键单击编辑区中的 \code{main.py}，选择运行该文件，并在运行窗口里核对输出。

\begin{figure}[htbp]
  \centering
  \begin{tikzpicture}
    \node[python flow step,text width=25mm] (python) {安装\\Python 3};
    \node[python flow step,text width=25mm,right=5mm of python] (venv)
      {创建项目\\与 \code{.venv}};
    \node[python flow step,text width=25mm,right=5mm of venv] (interpreter)
      {选择项目\\解释器};
    \node[python flow step,text width=25mm,right=5mm of interpreter] (run)
      {运行\\\code{main.py}};
    \draw[python flow arrow] (python) -- (venv);
    \draw[python flow arrow] (venv) -- (interpreter);
    \draw[python flow arrow] (interpreter) -- (run);
  \end{tikzpicture}
  \caption*{创建虚拟环境、选择解释器并运行脚本的配置流程。}
\end{figure}

\begin{pycharmpractice}
碰到导入失败时，先用下面的脚本打印 \code{sys.executable}。
在 PyCharm 和终端各跑一次，再比较两边的解释器路径。
\end{pycharmpractice}

\begin{lstlisting}[title={\code{check_interpreter.py}}]
import sys

# 打印当前进程所用的 Python 解释器路径。
print(sys.executable)
\end{lstlisting}

\begin{pycharmtip}
想少挪一次鼠标时，可双击 \code{Shift} 打开全局搜索，输入 \code{Run} 找到运行操作。
如果改过键位，就到 \code{Keymap} 查实际绑定。
\end{pycharmtip}

运行前，可按图~\ref{fig:pycharm-interpreter} 的关系核对解释器路径。
\begin{figure}[htbp]
  \centering
  \BookImage{interpreter}{\linewidth}
  \caption{项目与终端选用同一虚拟环境，解释器路径应与 \code{sys.executable} 一致。}
  \label{fig:pycharm-interpreter}
\end{figure}


\section{第一眼先看报错的哪一行}

将回溯当成一张倒着读的路线单。先看末尾的异常类型和错误消息，再往上找最近一处
属于自己脚本的文件名和行号。这样可以直接回到当时出问题的那次调用。

\begin{pitfall}
从控制台抄代码时，最容易顺手把 \code{>>>} 也带上。把
\code{>>> print(2 + 3)} 整行贴进脚本，会得到 \code{SyntaxError}。
保存到 \code{.py} 文件时，只留右侧的 \code{print(2 + 3)}。
\end{pitfall}

\begin{pybug}
将调用改成 \code{add("2", 3)} 时，错误会在函数里的加法处暴露出来。
注解给了一个预期，实际传进来的仍是一个字符串和一个整数。
如果这个字符串来自输入框，先在读取输入的地方用 \code{int()} 转换，
再把整数交给 \code{add()}。
\end{pybug}

图~\ref{fig:pycharm-traceback} 把异常末行与代码位置连起来，方便沿回溯走回现场。
\begin{figure}[htbp]
  \centering
  \BookImage{traceback}{\linewidth}
  \caption{\code{add("2", 3)} 在加法处触发 \code{TypeError}，回溯记录函数位置与异常类型。}
  \label{fig:pycharm-traceback}
\end{figure}

\subsection*{两分钟动手演练}

先把 \code{add(2, 3)} 改成 \code{add(-2, 3)}，在运行前猜一下结果。
接着恢复原例，再故意传入字符串 \code{"2"}，按上面的路线读一遍回溯。
最后改成 \code{add(int("2"), 3)}，看看程序又打印了什么。

\subsection*{快速起步的验收}

在 \code{return a + b} 这一行设断点，确认参数值是 \code{2} 和 \code{3}，
再继续运行，检查输出是否为 \code{5}。如果运行窗口和断点里的值对不上，
先检查当前运行配置指向哪个脚本，再检查解释器路径。

\subsection*{进阶边界演练}
可复制一份运行配置，把脚本换成另一个同名文件，再观察解释器路径和文件路径。这个练习说明，编辑器里打开的那份文件，和运行按钮真正启动的文件，需要用现场信息核对。
